% TV4 phu trach -- Thuc nghiem va Ket qua

% =====================================================================
\section{Thực nghiệm và Kết quả}
% =====================================================================

Phần thực nghiệm của nhóm được xây dựng theo \textbf{Cách 2: tự xây dựng demo từ đầu}.
Mục tiêu không phải tái hiện đầy đủ toàn bộ benchmark của bài báo Diffusion-DICE, mà là
tự cài đặt một phiên bản rút gọn để minh họa cơ chế trung tâm trong tutorial: khi một
generative policy dùng learned critic để guide hoặc select action, critic có thể bị khai thác
ngoài vùng dữ liệu; cơ chế in-sample guidance của Diffusion-DICE giúp giảm rủi ro này.

Khác với bản toy case ban đầu, demo chính hiện tại \textbf{không sử dụng dữ liệu tổng hợp}.
Nhóm dùng hai real benchmark datasets có sẵn trong scikit-learn: \textbf{Breast Cancer
Wisconsin Diagnostic} và \textbf{Wine Recognition}. Toàn bộ demo được chạy trong notebook
\texttt{demo/real\_benchmark\_diffusion\_dice\_demo.ipynb}; mã nguồn backend nằm trong
\texttt{demo/real\_benchmark\_diffusion\_dice.py}.

% =====================================================================
\subsection{Thiết lập bài toán}
% =====================================================================

Dữ liệu gốc là dữ liệu bảng nhiều chiều. Để có thể trực quan hóa quá trình sinh hành động,
nhóm chuẩn hóa dữ liệu rồi chiếu xuống không gian hai chiều bằng PCA. Mỗi điểm dữ liệu sau
PCA được xem như một action $a \in \mathbb{R}^2$ trong một bài toán contextual bandit rút
gọn. Nhãn thật của dataset được dùng để xác định vùng target class, tức vùng có reward cao.

Reward dùng để đánh giá không được đặt thủ công theo hình dạng tổng hợp. Thay vào đó, nhóm
xây dựng một \textbf{KNN oracle reward} từ nhãn thật:
\[
    r(a) =
    \frac{\sum_i K(a,a_i)\mathbf{1}[y_i = y^\star]}
         {\sum_i K(a,a_i)} ,
\]
trong đó $a$ là action được sinh ra, $a_i$ là điểm dữ liệu thật, $y_i$ là nhãn thật,
$y^\star$ là target class, và $K(a,a_i)$ là kernel đo độ gần giữa action và dữ liệu thật.
Vì reward này được suy ra từ nhãn thật, nó được dùng như thước đo khách quan hơn so với
điểm của learned critic.

Hai dataset được dùng như sau:
\begin{itemize}
    \item \textbf{Breast Cancer Wisconsin Diagnostic}: target class là \texttt{benign}.
    Sau PCA, hai thành phần đầu giải thích khoảng $63.2\%$ phương sai.
    \item \textbf{Wine Recognition}: target class là \texttt{class\_0}. Sau PCA, hai
    thành phần đầu giải thích khoảng $55.0\%$ phương sai.
\end{itemize}

% =====================================================================
\subsection{Phương pháp và baseline}
% =====================================================================

Nhóm tự cài đặt ba hướng sinh action trên cùng một dữ liệu offline:
\begin{enumerate}
    \item \textbf{IDQL-style select-only}: học behavior diffusion model từ dữ liệu thật,
    sinh nhiều candidate actions, sau đó chọn các action có learned critic score cao nhất.
    Rủi ro của hướng này là critic có thể đánh giá quá cao các candidate nằm xa data support.

    \item \textbf{QGPO-style guide-only}: dùng gradient của learned critic để guide quá
    trình reverse diffusion:
    \[
        a \leftarrow a + \eta \nabla_a Q_\phi(a).
    \]
    Ở đây $Q_\phi(a)$ là learned critic và $\eta$ là guidance scale. Nếu critic không đáng
    tin ở vùng ít dữ liệu, gradient này có thể kéo action ra khỏi phân phối thật.

    \item \textbf{Diffusion-DICE-style guide-then-select}: học DICE weights trên dữ liệu
    offline, sau đó dùng in-sample guidance để kéo candidate về vùng dữ liệu thật có trọng
    số cao:
    \[
        g(a) = \log \sum_i w_i K(a,a_i).
    \]
    Trong đó $w_i$ là trọng số DICE của điểm dữ liệu thật $a_i$. Sau bước guide, nhóm vẫn
    dùng learned critic để select, nhưng chỉ trong tập candidate đã được giữ gần data support.
\end{enumerate}

Các baseline trong demo là phiên bản \textit{style/proxy} của IDQL và QGPO nhằm minh họa
cơ chế guide-only và select-only. Nhóm không claim đây là full implementation của hai paper
IDQL hoặc QGPO.

% =====================================================================
\subsection{Metric đánh giá}
% =====================================================================

Mỗi phương pháp được đánh giá bằng cả metric định lượng và hình ảnh trực quan:
\begin{itemize}
    \item \textbf{ID ratio}: tỉ lệ generated actions nằm gần support của dữ liệu thật.
    \item \textbf{OOD ratio}: tỉ lệ generated actions nằm ngoài support.
    \item \textbf{Oracle reward mean}: reward trung bình theo KNN oracle từ nhãn thật.
    \item \textbf{Critic reward mean}: score trung bình theo learned critic.
    \item \textbf{Target region ratio}: tỉ lệ generated actions nằm trong vùng target class.
    \item \textbf{Mean nearest distance}: khoảng cách trung bình từ generated actions tới
    điểm dữ liệu thật gần nhất.
\end{itemize}

Nếu critic reward cao nhưng ID ratio thấp, phương pháp có dấu hiệu khai thác lỗi critic
(\textit{critic exploitation}). Ngược lại, một phương pháp tốt cần vừa giữ ID ratio cao,
vừa đạt oracle reward cao.

% =====================================================================
\subsection{Kết quả trên Breast Cancer}
% =====================================================================

\begin{figure}[H]
    \centering
    \includegraphics[width=1.0\textwidth]{images/real_breast_cancer_results_notebook.png}
    \caption{Kết quả trên Breast Cancer Wisconsin Diagnostic. Hàng trên lần lượt là data
    support sau PCA, KNN oracle reward từ nhãn thật, và learned critic score. Hàng dưới là
    action được sinh bởi QGPO-style, IDQL-style và Diffusion-DICE.}
    \label{fig:real_breast_cancer_results}
\end{figure}

\begin{table}[H]
    \centering
    \caption{Metric comparison trên Breast Cancer Wisconsin Diagnostic.}
    \label{tab:real_breast_cancer_metrics}
    \begin{tabular}{lccc}
        \hline
        \textbf{Method} & \textbf{ID ratio} & \textbf{OOD ratio} & \textbf{Oracle reward} \\
        \hline
        QGPO-style guide-only & 0.953 & 0.047 & 0.964 \\
        IDQL-style select-only & 1.000 & 0.000 & 0.961 \\
        Diffusion-DICE & 1.000 & 0.000 & 0.947 \\
        \hline
    \end{tabular}
\end{table}

Trên Breast Cancer, cả IDQL-style và Diffusion-DICE đều đạt reward cao và giữ action trong
vùng dữ liệu thật. QGPO-style cũng đạt oracle reward cao, nhưng đã có một phần action rơi ra
ngoài support dữ liệu (OOD ratio $=0.047$). Điểm đáng chú ý là Diffusion-DICE có mean nearest
distance thấp nhất trong ba phương pháp, nghĩa là action sinh ra bám sát dữ liệu thật hơn.

% =====================================================================
\subsection{Kết quả trên Wine Recognition}
% =====================================================================

\begin{figure}[H]
    \centering
    \includegraphics[width=1.0\textwidth]{images/real_wine_results_notebook.png}
    \caption{Kết quả trên Wine Recognition. Các panel gồm data support, oracle reward,
    learned critic, và action sinh bởi ba phương pháp.}
    \label{fig:real_wine_results}
\end{figure}

\begin{table}[H]
    \centering
    \caption{Metric comparison trên Wine Recognition.}
    \label{tab:real_wine_metrics}
    \begin{tabular}{lccc}
        \hline
        \textbf{Method} & \textbf{ID ratio} & \textbf{OOD ratio} & \textbf{Oracle reward} \\
        \hline
        QGPO-style guide-only & 0.984 & 0.016 & 0.750 \\
        IDQL-style select-only & 0.984 & 0.016 & 0.931 \\
        Diffusion-DICE & 1.000 & 0.000 & 0.916 \\
        \hline
    \end{tabular}
\end{table}

Trên Wine Recognition, IDQL-style đạt oracle reward cao nhất trong lần chạy này, cho thấy
baseline select-only vẫn rất mạnh trong setting nhỏ và đơn giản. Tuy nhiên, Diffusion-DICE
là phương pháp duy nhất giữ OOD ratio bằng 0 và ID ratio bằng 1.0. Điều này phù hợp với vai
trò chính của Diffusion-DICE: không chỉ tối đa hóa critic score, mà còn kiểm soát generated
actions để không rời khỏi support của dữ liệu thật.

% =====================================================================
\subsection{Mở rộng: khảo sát guidance scale}
% =====================================================================

Để đáp ứng yêu cầu mở rộng thực nghiệm, nhóm khảo sát guidance scale:
\[
    \eta \in \{0, 0.5, 1.0, 2.0, 3.0, 4.0\}.
\]
QGPO-style dùng $\eta$ để khuếch đại gradient của learned critic, trong khi Diffusion-DICE
dùng $\eta$ trên guidance có ràng buộc in-sample. Hai metric chính trong khảo sát là ID ratio
và oracle reward mean.

\begin{figure}[H]
    \centering
    \includegraphics[width=0.95\textwidth]{images/real_breast_cancer_tuning_notebook.png}
    \caption{Guidance-scale sweep trên Breast Cancer. Khi $\eta$ tăng lớn, QGPO-style mất
    dần ID ratio, trong khi Diffusion-DICE duy trì ID ratio bằng 1.0.}
    \label{fig:real_breast_cancer_tuning}
\end{figure}

\begin{figure}[H]
    \centering
    \includegraphics[width=0.95\textwidth]{images/real_wine_tuning_notebook.png}
    \caption{Guidance-scale sweep trên Wine Recognition. Diffusion-DICE giữ ID ratio ổn định
    và duy trì oracle reward cao trên nhiều giá trị $\eta$.}
    \label{fig:real_wine_tuning}
\end{figure}

Kết quả sweep cho thấy khác biệt chính giữa critic guidance và in-sample guidance. Với
Breast Cancer, QGPO-style có oracle reward tăng khi $\eta$ lớn, nhưng ID ratio giảm mạnh,
cho thấy một phần action đã bị kéo ra khỏi support dữ liệu thật. Diffusion-DICE giữ ID ratio
bằng 1.0 trên toàn bộ dải $\eta$. Với Wine Recognition, Diffusion-DICE cũng giữ ID ratio bằng
1.0 và oracle reward ổn định hơn.

% =====================================================================
\subsection{Mô tả demo và khả năng tái lập}
\label{sec:demo_description}
% =====================================================================

Demo Cách 2 được tổ chức trong thư mục \texttt{demo/} với các file chính:
\begin{itemize}
    \item \texttt{demo/real\_benchmark\_diffusion\_dice\_demo.ipynb}: notebook chính,
    dùng để chạy từng bước, hiển thị hình và metric.
    \item \texttt{demo/real\_benchmark\_diffusion\_dice.py}: backend tự implement
    learned critic, behavior diffusion, DICE weights, in-sample guidance, baseline,
    metric và visualization.
    \item \texttt{demo/requirements.txt}: danh sách thư viện cần cài đặt.
    \item \texttt{demo/real\_breast\_cancer\_*}: hình và CSV metric cho Breast Cancer.
    \item \texttt{demo/real\_wine\_*}: hình và CSV metric cho Wine Recognition.
\end{itemize}

Để chạy lại demo:
\begin{verbatim}
python -m pip install -r demo/requirements.txt
jupyter notebook demo/real_benchmark_diffusion_dice_demo.ipynb
\end{verbatim}
Hoặc chạy bằng script:
\begin{verbatim}
python demo/real_benchmark_diffusion_dice.py \
    --dataset all \
    --experiment all
\end{verbatim}

Demo chạy bằng Python và PyTorch, không cần GPU hay simulator như MuJoCo. Các output chính
được lưu dưới dạng PNG và CSV, giúp người đọc kiểm tra cả kết quả định tính lẫn định lượng.

% =====================================================================
\subsection{Hạn chế}
% =====================================================================

Demo này là một phiên bản rút gọn để minh họa cơ chế của Diffusion-DICE, không phải
reproduction đầy đủ của bài báo trên D4RL. PCA 2D giúp trực quan hóa rõ ràng nhưng làm mất
một phần thông tin từ dữ liệu gốc. Ngoài ra, các baseline IDQL-style và QGPO-style là proxy
để minh họa hai hướng select-only và guide-only, không phải full implementation của các paper
tương ứng. Vì vậy, kết luận chính của demo là về cơ chế kiểm soát phân phối và giảm critic
exploitation, không phải so sánh state-of-the-art trên benchmark lớn.
